% part: intuitionistic-logic
% chapter: soundness-completeness
% section: canonical-model
%READERNOTE{SOL6-A201: मूळ one-based pair list आणि शून्यासहित natural-number worlds यांतील विसंगती दुरुस्त केली. सर्व नैसर्गिक निर्देशांकांसाठी formula-pair enumeration शून्यापासून सुरू केली; शाखा आणि truth-lemma साठी लागणाऱ्या गुणधर्मांना बदल नाही.}

\documentclass[../../../include/open-logic-section]{subfiles}

\begin{document}

\olfileid{int}{sc}{mod}

\olsection{कॅनॉनिकल प्रतिमान}

आपल्या प्रतिमानातील जगे नैसर्गिक संख्यांच्या
सांत क्रमिका~$\sigma$ असतील, म्हणजे
$\sigma \in \Nat^*$. $\Nat^*$ ची
विगमनाने पुढीलप्रमाणे व्याख्या होते:
\begin{enumerate}
\item $\emptyseq \in \Nat^*$.
\item $\sigma \in \Nat^*$ आणि $n \in \Nat$
  असतील, तर $\sigma.n \in
  \Nat^*$ (येथे $\sigma.n$ म्हणजे
  $\sigma \concat \tuple{n}$, आणि
  $\sigma \concat \sigma'$ म्हणजे
  $\sigma$ व $\sigma'$ यांची जोडणी).
\item इतर काहीही $ \Nat^*$ मध्ये नाही.
\end{enumerate}
म्हणून $\Nat^*$ वापरून विगमनात्मक
व्याख्या देता येतात.

!!{formula}s च्या सर्व जोड्यांचे
$\tuple{!B_0, !C_0}$,
$\tuple{!B_1, !C_1}$,
$\tuple{!B_2, !C_2}$, \dots,
असे $n\in\Nat=\{0,1,2,\dots\}$ ने निर्देशांकित
प्रगणन असू द्या. !!{formula}s
चा संच~$\Delta$ दिला असता
$\Delta(\sigma)$ ची विगमनाने
पुढीलप्रमाणे व्याख्या करा:
\begin{enumerate}
\item $\Delta(\emptyseq) = \Delta$
\item $\Delta(\sigma.n) = {}$
  \[
  \begin{cases}
    (\Delta(\sigma) \cup \{!B_n\})^* &
    \text{जर $\Delta(\sigma) \cup \{!B_n\} \Proves/ !C_n$ असेल} \\
    \Delta(\sigma) & \text{अन्यथा}
  \end{cases}
  \]
\end{enumerate}
येथे $(\Delta(\sigma) \cup \{!B_n\})^*$
म्हणजे $\Delta(\sigma) \cup \{!B_n\}$
या संचाला आणि !!{formula}~$!C_n$
ला \olref[lin]{lem:lindenbaum} लावल्यावर
मिळणारा !!{formula}s चा अभाज्य संच.
या व्याख्येनुसार
$\Delta(\sigma) \cup \{!B_n\} \Proves/ !C_n$
असेल, तर $\Delta(\sigma.n)
\Proves !B_n$ आणि
$\Delta(\sigma.n) \Proves/ !C_n$,
हे लक्षात घ्या. तसेच कोणत्याही~$n$
साठी $\Delta(\sigma) \subseteq \Delta(\sigma.n)$.
$\Delta$ अभाज्य असेल, तर सर्व~$\sigma$
साठी $\Delta(\sigma)$ अभाज्य असतो.

\begin{defn}\ollabel{defn:canonical-model}
  $\Delta$ अभाज्य आहे असे समजा.
  मग $\Delta$ साठीचे
  \emph{कॅनॉनिकल प्रतिमान}
  $\mModel{M(\Delta)}$ पुढीलप्रमाणे परिभाषित होते:
  \begin{enumerate}
  \item $W = \Nat^*$, म्हणजे नैसर्गिक
    संख्यांच्या सांत क्रमिकांचा संच.
  \item $R$ हा असा अंशतः क्रम आहे की
    $R\sigma\sigma'$ तेव्हा आणि केवळ
    तेव्हाच जेव्हा $\sigma$ हा~$\sigma'$
    चा आरंभीचा खंड असतो (म्हणजे
    काही क्रमिका~$\sigma''$ साठी
    $\sigma' = \sigma \concat \sigma''$).
  \item $V(p) = \Setabs{\sigma}{p \in \Delta(\sigma)}$.
  \end{enumerate}
\end{defn}

$R$ हा खरोखर अंशतः क्रम आहे,
हे सहज तपासता येते. तसेच~$V$ वरील
एकस्वनिकतेची अट पूर्ण होते.
$\Delta(\sigma) \subseteq \Delta(\sigma.n)$
असल्याने $R\sigma\sigma'$ असेल तेव्हा
$\Delta(\sigma) \subseteq \Delta(\sigma')$;
हे~$\sigma$ पासूनच्या विस्ताराच्या
लांबीवर विगमनाने मिळते.

\end{document}
